\begin{figure}[htbp]
\centering
\begin{tikzpicture}[
node distance=8mm and 6mm,
box/.style={draw, rounded corners, align=center, minimum height=8mm, text width=2.0cm},
arrow/.style={-{Latex[length=2mm]}, thick}
]
\node[box] (challenge) {Challenge};
\node[box, right=of challenge] (destab) {Destabili-\\zation};
\node[box, right=of destab] (pause) {Pause};
\node[box, right=of pause] (orient) {Orient};
\node[box, below=of orient] (strategy) {Select strategy};
\node[box, left=of strategy] (regulate) {Regulate};
\node[box, left=of regulate] (recover) {Recover};
\node[box, left=of recover] (resume) {Resume goal};

\draw[arrow] (challenge) -- (destab);
\draw[arrow] (destab) -- (pause);
\draw[arrow] (pause) -- (orient);
\draw[arrow] (orient) -- (strategy);
\draw[arrow] (strategy) -- (regulate);
\draw[arrow] (regulate) -- (recover);
\draw[arrow] (recover) -- (resume);
\draw[arrow] (resume.west) |- ++(-5mm,-8mm) -| (challenge.south);
\end{tikzpicture}
\caption{Core ARC-ADHD regulation loop. A successful episode may contain destabilization or error; the target is recovery and return to the interrupted goal.}
\label{fig:closed-loop}
\end{figure}
